\subsubsection{Dependent Function Types ($\Pi$-types)}
There is an extension of function types, of `dependent functions' where the codomain type may depend
on the doamin. Given $A : \U, B : A \to \U$, the type of dependent functions is
\begin{align*}
  \Prod{x:A} B(x)
\end{align*}
which now given $a : A$, give $f(a) : B(a)$ for any instance $f : \Prod{x:A} B(x)$.

\subsubsection{Product Types}
Given $A, B : \U$, the product type $A \times B : \U$ consists of pairs $(a, b)$, constructed by $\mathsf{pair} : A \to (B \to A \times B)
To$ 

\subsubsection{Coproduct Types}

\subsubsection{Dependent Pair Types ($\Sigma$-types)}

\subsubsection{0, 1 and 2}
The zero type $\0$ is the type with no constructors, and the induction principle, 
\begin{align*}
  \ind_\0 : \Prod{C : \0 \to \U}{ \Prod{x : \0} {C(x)}}
\end{align*}
with no computation principle since there is no constructor.

The unit type $\1$ is the type with a single term as a constructor, $\star : 1$, and the induction principle
\begin{align*}
  \ind_{\1} : \Prod{C : \1 \to \U}{C(\star) \to \left(\Prod{x : 1} C(x)\right)}
\end{align*}
The boolean type $\2$ is defined by two constructors $0_2 : \2$, and $1_2 : \2$, or equivalently as $1 + 1$.

\subsubsection{Natural Numbers Type ($\N$)}

\subsubsection{Identity Types}

\textbf{Remark.} Inductive Types may be interpreted as ...
freely generated by constructors.
Lean, Coq, Agda


\begin{proposition}
  There exist maps $(\cdot)^{-1} : \Prod{x, y : A}{(x =_A y) \to (y =_A x)}$ and $(-\sqcdot-) : \Prod{x,y,z :A}{(x =_A y) \to (y =_A z) \to (x =_A z)}$
  such that for $x, y, z, w : A$, $p : x =_A y, q : y =_A z$ and $r : z =_A w$, we have
\end{proposition}

\begin{proposition}
  Given $x, y, z : A$, there exist the following maps 
  \begin{itemize}
    \item [(i)] $(x = y) \to (y = x)$
    \item [(ii)] $(x = y) \to (y = z) \to (x = z)$
  \end{itemize}
\end{proposition}
\textit{Proof.} (i) By path induction, assume $x \deq y$ and $p \deq \refl_x$. Now simply exhibit $\mathsf{id}_{x =_A x} \deq (\lambda a.a)$. 
(ii) By path induction, assume $x \deq y$ and $p \deq \refl_x$. Exhibit $(\lambda a.\lambda b.b)$. \qed

Denote the map in $(i)$ as $(-)^{-1}$, and $(ii)$ as $(- \sqcdot -)$. 

\begin{proposition} Let $x, y, z, w : A$, $p : x =_A y, q : y =_A z$ and $r : z =_A w$. Then we have:
  \begin{itemize}
    \item [(i)] $p = \refl_x \sqcdot p$, and $p = p \sqcdot \refl_y$
    \item [(ii)] $p \sqcdot p^{-1} = \refl_x$ and $p^{-1} \sqcdot p = \refl_y$
    \item [(iii)] $(p \sqcdot q) \sqcdot r = p \sqcdot (q \sqcdot r)$ 
   \end{itemize} 
\end{proposition}
\textit{Proof.} (i) By path induction, assume $x \deq y$ and $p \deq \refl_x$. Now exhibit $\refl_{\refl_x}$ or $\refl_{\refl_y}$.
(ii) By path induction, assume $x \deq y$ and $p \deq \refl_x$. Exhibit $\refl_{\refl_x}$ or $\refl_{\refl_y}$.
(iii) By path induction, assume $y \deq z$ and $q \deq \refl_y$. Exhibit $\refl_p$ to obtain $(p \sqcdot q) \deq  p$,
and similarly $\refl_r$ to obtain $(q \sqcdot r) \deq r$. Our goal reduces to $p \sqcdot r = p \sqcdot r$.
Exhibit $\refl_{p \sqcdot r} : p \sqcdot r = p \sqcdot r$.
\qed

A remarkable feature is that while (i) and (ii) can be proven judgementally, (iii) is only true propositionally,
i.e., upto a path in the higher path space.
Strikingly, we also have all the required higher coherences propositionally using only path induction.
This gives us the (weak) $\infty$-groupoid interpretation which is made functorial by the following result.

\begin{proposition}
Let $f : A \to B, g : \Prod{a:A}{P(a)}, x, y : A$, and $p : x =_A y$, then we have the following: 
  \begin{itemize}
    \item [(i)] (Application) $\mathsf{ap}_f(p) : (f(x) =_B f(y))$
    \item [(ii)] (Transport) $p_* : P(x) \to P(y)$ 
    \item [(iii)] (Dependent Application) \\ $\mathsf{apd}_g(p) : (p_*(g(x)) =_{P(y)} g(y))$ 
  \end{itemize}
\end{proposition}
\textit{Proof.} For all, use path induction. Assume $x \deq y$ and $p \deq \refl_x$ for all.
(i) Exhibit $\refl_{f(x)}$.
(ii) Exhibit $\mathsf{id}_{P(x)} :\deq (\lambda z.z)$.
(iii) Exhibit $\refl_{g(x)}$.
\qed


